\DocumentMetadata{
  pdfversion=2.0,
  pdfstandard=ua-2,
  lang=en-US,
  tagging=on,
  tagging-setup={math/setup=mathml-SE,tabsorder=structure}
}
\documentclass[11pt,letterpaper]{article}
\usepackage[margin=0.68in,includefoot]{geometry}
\usepackage{newcomputermodern}
\usepackage{amsmath,amscd,mathtools}
\providecommand{\square}{\mdlgwhtsquare}
\usepackage[hidelinks]{hyperref}
\hypersetup{
  pdftitle={Analysis First-Year Exam, August 2025, Part 1},
  pdfauthor={University of Florida Department of Mathematics},
  pdfsubject={Graduate mathematics examination},
  pdfdisplaydoctitle=true,
  bookmarksopen=true
}
\setcounter{secnumdepth}{0}
\setlength{\parindent}{0pt}
\setlength{\parskip}{0.28em}
\setlength{\emergencystretch}{3em}
\setlength{\leftmargini}{2.15em}
\setlength{\leftmarginii}{2.65em}
\setlength{\leftmarginiii}{3em}
\renewcommand{\labelenumi}{\textbf{\theenumi.}}
\pagestyle{empty}
\raggedbottom
\allowdisplaybreaks
\newenvironment{examproblems}[1]{%
  \begin{enumerate}%
  \setcounter{enumi}{#1}%
  \setlength{\topsep}{0.35em}%
  \setlength{\partopsep}{0pt}%
  \setlength{\itemsep}{0.48em}%
  \setlength{\parsep}{0.18em}%
}{\end{enumerate}}
\newenvironment{examsubparts}{%
  \begin{enumerate}%
  \setlength{\topsep}{0.18em}%
  \setlength{\partopsep}{0pt}%
  \setlength{\itemsep}{0.16em}%
  \setlength{\parsep}{0.1em}%
}{\end{enumerate}}
\newenvironment{examinstructions}{%
  \par\begingroup\itshape
}{%
  \par\endgroup\vspace{0.18em}
}
\makeatletter
\renewcommand\section{\@startsection{section}{1}{0pt}%
  {-1.25ex plus -.25ex minus -.1ex}%
  {0.55ex plus .1ex}%
  {\normalfont\large\bfseries}}
\renewcommand{\@maketitle}{%
  \newpage\null\vskip -1.2em%
  {\footnotesize\bfseries
    \href{https://www.ufl.edu/}{University of Florida}, \href{https://math.ufl.edu/}{Department of Mathematics}\par}%
  \vskip 0.18em%
  \tagstructbegin{tag=Title}%
    {\LARGE\bfseries\href{https://gma.math.ufl.edu/exams/analysis-first-year/}{Analysis First-Year Exam}, August 2025, Part 1\par}%
  \tagstructend%
  \vskip 0.35em%
  \tagmcbegin{artifact}%
    \hrule height 1pt%
  \tagmcend%
  \vskip 0.55em%
}
\makeatother
\title{Analysis First-Year Exam, August 2025, Part 1}
\author{}
\date{}
\begin{document}
\maketitle
\thispagestyle{empty}
\begin{examinstructions}
Answer FOUR questions in detail. State carefully any results used without proof.
\end{examinstructions}
\begin{examproblems}{0}
\item
Let \(A \subseteq \mathbb{R}\) be closed under subtraction, in the sense that if \(s,t \in A\) then \(s-t \in A\). Prove that if \(A \cap (0,u) \ne \varnothing\) whenever \(u>0\) then~\(A\) is dense in~\(\mathbb{R}\).
\item
For every positive integer~\(n\) let \(A_n\) be a subset of the metric space~\(X\). For each of the following statements, give a proof or a counterexample. Here, an overline signifies closure.
\begin{examsubparts}
\item[\textnormal{(i)}]
if~\(N\) is a positive integer then \(\overline{\bigcup_{n=1}^N A_n}=\bigcup_{n=1}^N \overline{A_n}\);
\item[\textnormal{(ii)}]
\(\overline{\bigcup_{n=1}^\infty A_n}=\bigcup_{n=1}^\infty \overline{A_n}\).
\end{examsubparts}
\item
Let~\(S\) be a countably infinite subset of \(\mathbb{R}^2\) (with the Euclidean metric). Must the complement \(\mathbb{R}^2 \setminus S\) be connected? If so, give a proof; if no, give a counterexample.
\item
Let the nonempty closed subsets~\(A\) and~\(B\) of a metric space be disjoint. Must the infimum \(\inf\{d(a,b):a\in A,b\in B\}\) be strictly positive? Does the answer change if either~\(A\) or~\(B\) is compact?
\item
Let \(f:\mathbb{R}\to\mathbb{R}\) be differentiable, with \(f'(0)=1\). Must there exist an open interval about~\(0\) on which~\(f\) is increasing? Proof or counterexample.
\end{examproblems}
\end{document}
